\chapter{FUNDAMENTAÇÃO TEÓRICA}

Considerando o caráter prático deste trabalho, a fundamentação teórica foi construída a partir de uma abordagem multidisciplinar, reunindo conceitos relacionados aos sistemas de informação, à transformação digital, à automação de processos e à arquitetura de software para aplicações web. Esses temas fornecem a base conceitual necessária para o desenvolvimento de um Sistema Integrado de Gestão (ERP) direcionado às marmorarias, um setor que ainda apresenta forte dependência de processos manuais, controles descentralizados e baixo nível de digitalização.

A base teórica adotada contempla tanto os aspectos tecnológicos quanto os impactos organizacionais decorrentes da implementação de soluções digitais em pequenas e médias empresas. Dessa forma, busca-se compreender como a utilização de sistemas integrados pode contribuir para a otimização das atividades operacionais, a redução de falhas e retrabalhos, o aumento da rastreabilidade das informações e o suporte à tomada de decisões estratégicas.

Ao longo deste capítulo são apresentados os principais conceitos que fundamentam a proposta do sistema, incluindo sistemas de informação, transformação digital, automação de processos, integração entre setores, usabilidade, computação em nuvem e práticas modernas de desenvolvimento de software. Esses elementos constituem a base para compreender como tecnologias atuais podem estruturar e padronizar processos que, tradicionalmente, são executados de forma manual e pouco integrada.

Assim, a fundamentação teórica estabelece a relação entre os desafios enfrentados pelas marmorarias e os conceitos que orientam o desenvolvimento da solução proposta. A partir dessa perspectiva, os tópicos seguintes apresentam os referenciais que sustentam a concepção e a implementação do ERP, evidenciando como a tecnologia pode promover maior eficiência, organização e competitividade para empresas do setor.


\section{A Gestão Tradicional versus a Transformação Digital nas Marmorarias}
A gestão operacional é um dos pilares para o funcionamento eficiente de empresas do setor de rochas ornamentais. No contexto de uma marmoraria, tarefas como orçamento, medição, produção, controle de pedidos, comunicação com clientes e instalação demandam precisão e organização. Entretanto, esses processos ainda são frequentemente conduzidos de forma manual, fragmentada e pouco padronizada, o que aumenta a ocorrência de falhas, perda de informações, atrasos e retrabalhos.

Com o avanço das Tecnologias da Informação e Comunicação (TICs), as empresas passaram a ter acesso a sistemas de gestão capazes de centralizar dados, automatizar tarefas e melhorar a comunicação interna. A transformação digital se tornou um fator decisivo para aumentar competitividade, especialmente em segmentos onde a tecnologia ainda é subutilizada. Isso inclui desde a modernização do fluxo operacional até o uso de ferramentas que possibilitam tomada de decisão baseada em dados.

Nesse cenário, um ERP especializado surge como resposta à necessidade de integrar processos, otimizar rotinas diárias e estruturar operações que antes dependiam de controles informais. A digitalização permite maior rastreabilidade, melhora a eficiência e reduz gargalos, além de criar um ecossistema onde colaboradores, gestores e clientes se beneficiam de informações claras e acessíveis.

A fundamentação teórica considera, portanto, que a adoção de plataformas digitais intuitivas, funcionais e alinhadas ao comportamento moderno dos usuários é essencial para que marmorarias alcancem níveis mais altos de produtividade. À medida que cresce a demanda por entregas rápidas e comunicação eficiente, torna-se indispensável que essas empresas acompanhem o ritmo da transformação digital global, incorporando sistemas que fortaleçam sua organização e crescimento sustentável.

\section{Sistemas Integrados de Gestão (ERP)}
Os Sistemas Integrados de Gestão (Enterprise Resource Planning – ERP) surgiram como uma solução tecnológica capaz de unificar informações, automatizar processos e centralizar operações antes distribuídas em diferentes setores de uma empresa. De acordo com a literatura, sistemas ERP desempenham papel essencial na padronização de rotinas, aumento da rastreabilidade operacional e suporte à tomada de decisão baseada em dados, promovendo maior eficiência e controle organizacional.

No contexto das marmorarias, a aplicação de sistemas ERP ainda é pouco difundida, embora o segmento apresente forte dependência de processos contínuos, como orçamentos, registros de clientes, medição, produção, logística e instalação. A ausência de integração entre essas etapas causa falhas de comunicação, atrasos e desperdícios, impactando diretamente a experiência do cliente e a lucratividade do negócio.

A adoção de um ERP especializado permite que essas operações, antes realizadas por meio de anotações, mensagens informais ou planilhas dispersas, sejam transformadas em um fluxo digital único e padronizado. Isso cria uma visão global do negócio, facilita o acompanhamento das etapas do serviço e possibilita um controle operacional mais preciso. Dessa forma, o ERP funciona como o núcleo organizacional que conecta setores, automatiza rotinas e cria um ambiente favorável à escalabilidade e ao crescimento sustentável.

\section{Computação em Nuvem e Arquitetura Serverless}
A computação em nuvem se consolidou como um dos pilares da modernização de sistemas corporativos, oferecendo infraestrutura escalável, alta disponibilidade e baixo custo operacional quando comparada a servidores próprios. A utilização de tecnologias como containers, pipelines de CI/CD e serviços automatizados permite que aplicações sejam desenvolvidas, distribuídas e mantidas com maior rapidez e eficiência.

A arquitetura serverless, utilizada em plataformas como o Google Cloud Run, representa um avanço significativo na forma como softwares são implantados e executados. Nesse modelo, o provedor de nuvem gerencia automaticamente a infraestrutura, escala conforme a demanda e cobra apenas pelo consumo efetivo de recursos. Para pequenas e médias empresas — como as marmorarias — essa abordagem elimina a necessidade de manter servidores dedicados, reduz custos e garante desempenho mesmo em picos de acesso.

Além disso, o uso de conteinerização e repositórios como o Artifact Registry permite padronizar ambientes, garantir reprodutibilidade e facilitar a gestão do ciclo de vida da aplicação. Aliado a pipelines de CI/CD, o desenvolvimento se torna mais ágil, seguro e confiável, possibilitando entregas contínuas e correções rápidas. Assim, a computação em nuvem, combinada à arquitetura serverless, proporciona a base tecnológica necessária para aplicações modernas, acessíveis e escaláveis — exatamente o que um ERP especializado necessita para operar com estabilidade e evolução constante.

\section{Boas Práticas de Desenvolvimento e Segurança}
A construção de um sistema de gestão exige mais do que funcionalidades operacionais; requer atenção às boas práticas de engenharia de software, garantindo que o produto final seja seguro, confiável e sustentável a longo prazo. Entre essas práticas estão a modularização de componentes, versionamento de código, testes automatizados, documentação clara e padronização da arquitetura da aplicação.

Um aspecto central é a segurança, especialmente em sistemas que lidam com dados sensíveis de clientes, pedidos e informações comerciais. A implementação de autenticação baseada em tokens JWT, por exemplo, permite controlar acessos de forma estruturada, garantindo que apenas usuários autorizados possam manipular dados internos do sistema. Esse mecanismo, aliado ao uso de protocolos seguros (HTTPS), gerenciamento adequado de permissões e armazenamento protegido, contribui para a integridade e a confidencialidade das informações.

Além disso, a adoção de boas práticas como monitoramento ativo, logs padronizados, versionamento semântico e ambientes separados para desenvolvimento, homologação e produção fortalece o ciclo de vida do software. Em conjunto, esses elementos garantem que o ERP seja não apenas funcional, mas também robusto, escalável e preparado para evoluções futuras, alinhando o projeto às exigências do desenvolvimento moderno.